\subsection{Backend mit FastAPI}
\label{subsec:fastapi-backend}

Das Backend ist eine separate HTTP-Schicht der Profile Web-cloud,
Desktop-cloud und Full-platform. \path{backend/app/main.py} erzeugt die
FastAPI-Anwendung, lädt typisierte Einstellungen, richtet CORS ein und bindet
den Router unter \texttt{/api} ein. Der derzeitige Router enthält ausschließlich
\texttt{/health} und \texttt{/ready}; produktspezifische Dienste, Fachmodelle
und Authentifizierung sind bewusst nicht vorgegeben. FastAPI stellt hierfür die
HTTP-Routen, Validierung und OpenAPI-Beschreibung bereit
\parencite{fastapi2026features,kleiveist2026architecture}.

Liveness meldet nur den Prozesszustand. Readiness prüft bei aktivierter
Datenbankfähigkeit zusätzlich mit \texttt{SELECT 1} und antwortet bei fehlender
Konfiguration oder Verbindung mit HTTP~503, ohne interne Fehler offenzulegen.
Tests adressieren Router, Einstellungen und diese Abhängigkeit getrennt. Die
Architektur hält Handler klein; erst ein abgeleitetes Produkt ergänzt
Service- oder Domänenmodule.

Auch im Deployment bleiben die Einheiten getrennt
(Abbildung~\ref{fig:deployment-architecture}). Der Vite-Build kann statisch
bereitgestellt werden, während FastAPI in einem eigenen, nicht privilegierten
Container läuft. Compose simuliert diese providerneutrale Grenze lokal;
PostgreSQL wird nur über ein explizites Profil ergänzt
\parencite{kleiveist2026deployment}.

\begin{figure}[H]
  \centering
  \resizebox{\textwidth}{!}{%
  \begin{tikzpicture}[node distance=8mm and 11mm]
    \node[casePrimary,minimum width=31mm] (frontend) {Frontend Build};
    \node[casePrimary,minimum width=31mm,right=of frontend] (backend) {Backend Service};
    \node[caseBox,minimum width=31mm,right=of backend] (container) {Docker / Compose};
    \node[caseBox,minimum width=26mm,below=of backend] (health) {Health};
    \node[caseBox,minimum width=26mm,right=of health] (readiness) {Readiness};
    \node[caseOptional,minimum width=36mm,right=of readiness] (postgres) {PostgreSQL optional};

    \draw[caseFlow] (frontend) -- (backend);
    \draw[caseFlow] (backend) -- (container);
    \draw[caseSecondaryFlow] (container) -- (health);
    \draw[caseSecondaryFlow] (container) -- (readiness);
    \draw[caseSecondaryFlow] (container) -- (postgres);
    \begin{scope}[on background layer]
      \node[caseGroup,fit=(frontend)(backend)(container)(health)(readiness)(postgres),
        label={[font=\sffamily\bfseries\scriptsize,text=caseBlue]above:providerneutrale Deployment-Grenze}] {};
    \end{scope}
  \end{tikzpicture}}
  \caption{Providerneutrale Deployment-Architektur}
  \label{fig:deployment-architecture}
\end{figure}

